\chapter{Conclusion and Outlook}\label{chap:sum}
%Our model of an \gls{ell}, or more specifically the \gls{sg}-like perturbation $H_2$ from \autoref{eqn:ell:bos:H2H4_model}, at finite temperature and for a non-trivial Luttinger parameter $K\neq 1$ demonstrably leads to the formation of \acrlongpl{ep}.
%From a high level point of view, it is therefore justified to deem the program of this thesis a success.
%Despite the unusual structure of the perturbation, one should not view this as a contrived academic example either; after all, the \gls{ell} corresponds to the bosonization of a concrete microscopic lattice model.\\
Starting from a concrete microscopic lattice model, bosonization leads us to a low-energy limit described by the \gls{ell} in \autoref{eqn:ell:bos:H2H4_model}.
The essential new component is the \gls{sg}-like perturbation $H_2$, which at finite temperature and for a non-trivial Luttinger parameter $K\neq 1$ demonstrably leads to the formation of \acrlongpl{ep}.
Qualitatively, the analytic results thus match the prediction from numerical data for the microscopic model in \cite{EP_DynamicInduced_Quenched_Lehmann_2021}.
The structural similarity of our \autoref{fig:ell:ep:ep_contour_dual_example} to Figure 2 in \cite{EP_DynamicInduced_Quenched_Lehmann_2021} is particularly striking.\\
Quantitative agreement can also be achieved, but at the time requires fitting model parameters to the data.
%As our \gls{rg} analysis demonstrates, it is conceivable that the corresponding values can also be predicted, but our perturbative one-loop calculation is insufficient for this task.
In principle, within our \gls{rg} analysis it is conceivable that the corresponding values can also be predicted, but our perturbative one-loop calculation is insufficient for this task.
%A possible avenue for future research would be to apply more sophisticated, non-perturbative renormalization schemes.
%This may provide an improvement over the perturbative approach, since we work far from the fixed point of the model.\\
Since we work far from the model's fixed point, one might gain from considering non-perturbative renormalization schemes, which provides an avenue for future research.\\
%However, it is also possible that quantitative agreement is hindered by one of the other encountered issues.
There are also other possibilities to pursue to try and improve the quantitative results.
%For one, the calculation of the single-particle \gls{gf} via perturbation theory is already quite difficult at first order, and, at least using our brute-force approach, seems to become completely infeasible beyond.
For one, using our brute-force approach, the calculation of the single-particle \gls{gf} via perturbation theory is only feasible to first order.
The situation could be improved by a better understanding of the integrals defining the functions $J_{b,N}$ from \autoref{eqn:app:J_bN_function}, because this might allow a higher order calculation.
As mentioned in a footnote in that section, the final result suggests a decomposition reminiscent of light-cone variables, but due to the complex nature of $\tau \pm \i x$, it is not yet clear how to exploit this.\\
%Of course abandoning the goal of exactly computing the integrals could also prove fruitful, but controlled approximations appear to be difficult.\\ %semi-classical path integral approach?
%Another problem is that we retain the \acrshort{uv} cut-off $\alpha$ for $K\neq 1$ in the final result, and although a finite value of $\alpha\propto\frac{1}{\beta v}$ can be argued via \gls{rg}, the analytic inconsistency with fundamental conditions on Matsubara \glspl{gf} remains.
Furthermore, for $K\neq 1$ we retain the \acrshort{uv} cut-off $\alpha$ in the final result.
Although a finite value of $\alpha\propto\frac{1}{\beta v}$ can be argued via \gls{rg}, the analytic inconsistency with fundamental conditions on Matsubara \glspl{gf} remains.
On a similar note, we can not quite take a proper continuum limit for the bosonized Hamiltonian, because setting $a=0$ eliminates the perturbation $H_2$.
Interestingly, this problem can not easily be solved by an expansion in the lattice spacing, as intermediate results are non-perturbative in $a$ due to the residue structure in certain integrals.
Thus, removing the mild non-locality induced by $a$ while retaining the \glspl{ep} may require some additional insight.\\
%Still, since we can not give a good argument to suggest the contrary, it should be possible to get rid of the mild non-locality induced by $a$ and still get \glspl{ep}.\\
Finally, the method used for the numerical results is limited to moderately high temperatures and interaction strengths.
%Both of these are less than ideal for a quantitative comparison to our low-order, field theoretic description, and in order to make progress it would be preferable to be able to at least analyse the asymptotic limit $U_\A,U_\B\rightarrow 0$.
Both of these make a quantitative comparison to our low-order, field theoretic description more challenging.
In order to make progress, it would be useful to be able to analyse the asymptotic limit $U_\A,U_\B\rightarrow 0$.
One approach might be to use plain perturbation theory directly in terms of the fermions to provide an analytic benchmark.
As discussed in \autoref{sec:ell:ep:quant_comp_num}, the special case of $U_\B=0$ is difficult to treat in our \gls{ell}.
The fermionic approach would work directly in the $\A\B$-basis, and therefore setting $U_\B=0$ should instead simplify the calculation, which provides another incentive for this idea.
%We hope that the problems can be resolved in the future and that our results spark interest in further research on \gls{nh} physics using bosonization and \glspl{ll}.
%The remaining problems may potentially be resolved in the future, and hopefully our results spark some interest in further research on \gls{nh} physics using bosonization and \glspl{ll}.